\section{What Is a Lattice}


\begin{definition}[Lattice]
A lattice $\Lambda$ in $\mathbb{R}^n$ is the set of all integer linear combinations of
a set of linearly independent vectors $\mathbf{b}_1, \ldots, \mathbf{b}_n \in \mathbb{R}^n$:
\[
\Lambda = \left\{ \sum_{i=1}^n z_i \mathbf{b}_i \;:\; z_i \in \mathbb{Z} \right\}
\]
The vectors $\{\mathbf{b}_i\}$ form a \emph{basis} for $\Lambda$.
\end{definition}

Think of a lattice as a grid---a regular arrangement of points in space. In two dimensions,
it looks like graph paper tilted and stretched. The crucial property: the grid is
\emph{discrete} (points are separated by gaps) but \emph{infinite} (it extends in all
directions).

\subsection{Why Lattices for Cryptography}

Classical public-key cryptography (RSA, elliptic curves) relies on the hardness of
factoring or discrete logarithms. Shor's algorithm solves both in polynomial time on a
quantum computer. Lattice problems---finding short vectors in high-dimensional
lattices---resist all known quantum algorithms.

The two fundamental hard problems:

\begin{definition}[Shortest Vector Problem (SVP)]
Given a lattice basis $\mathbf{B}$, find the shortest nonzero lattice vector.
\end{definition}

\begin{definition}[Closest Vector Problem (CVP)]
Given a lattice basis $\mathbf{B}$ and a target point $\mathbf{t}$, find the lattice
vector closest to $\mathbf{t}$.
\end{definition}

Both are NP-hard in their exact forms, and the best known algorithms (classical or
quantum) for approximate versions run in exponential time in the lattice dimension $n$.

\subsection{The Geometry of Hard Problems}

In low dimensions ($n = 2, 3$), finding short vectors is easy---look at the basis and
reduce it. In high dimensions ($n = 256, 512, 1024$), the geometry becomes pathological:
the shortest vector is exponentially shorter than the basis vectors, but finding it
requires exploring an exponentially large search space. This is the core intuition:
lattice cryptography works because high-dimensional geometry is hard.

%% ============================================================
